\pdfinfoomitdate=1
\pdftrailerid{}
\documentclass[10pt]{article}
\usepackage[a4paper,margin=0.55in]{geometry}
\usepackage[T1]{fontenc}
\usepackage{lmodern}
\usepackage{amsmath,amssymb,booktabs,microtype}
\linespread{0.96}
\title{A Finite Periodic Skeleton Disjointly Attached to an\\Aperiodic Minimal Subshift}
\author{Route-A structural certificate C149}
\date{}
\begin{document}
\maketitle
\small
\begin{abstract}
We attach tagged cycles of lengths $1,2,3,5$ by topological disjoint union to the aperiodic minimal
Thue--Morse subshift.  The resulting system has four primitive cycles and
rational Artin--Mazur zeta
$((1-z)(1-z^2)(1-z^3)(1-z^5))^{-1}$.  This supplies a controlled finite
periodic skeleton but destroys minimality.
The construction is a declared attachment, not an almost-minimal system or an
intrinsic creation of orbits inside the substitution language.
\end{abstract}

\section{The disjoint system}
Let $X_{\rm TM}$ be the two-sided subshift generated by
$0\mapsto01$, $1\mapsto10$.  It is nonempty and minimal and contains no
periodic point.  For $E=\{1,2,3,5\}$, let $C_\ell$ be a tagged cycle of length
$\ell$ and set
\[
Y=X_{\rm TM}\sqcup\bigsqcup_{\ell\in E}C_\ell .
\]
The map is the shift on $X_{\rm TM}$ and cyclic successor on each $C_\ell$.
One iterate is the clock.

\paragraph{Aperiodicity of the infinite component.}
For completeness, let $t_j$ be the parity of the binary digit sum of $j$.
If $0\leq r<2^q$, binary concatenation gives
$t_{j2^q+r}=t_j\mathbin{\mathsf{xor}}t_r$.  Fix a proposed period $p$ and
choose odd $k$ larger than its binary length.  For
\[
d=p(2^k-1)=(p-1)2^k+(2^k-p),
\]
the lower $k$ bits complement the $k$-bit expansion of $p-1$; hence
$\operatorname{popcount}(d)=k$ is odd.  Thus $p\mid d$ but $t_d\ne t_0$.
Let $b$ be the binary length of $d$.  Every interval of $t$ of length
$2^{b+1}$ contains a full block aligned at a multiple $j2^b$.  Its symbols at
offsets $0,d$ are $t_j\mathbin{\mathsf{xor}}t_0$ and
$t_j\mathbin{\mathsf{xor}}t_d$, hence are opposite at a distance divisible by
$p$.  No such interval is $p$-periodic.  A $p$-periodic point of $X_{\rm TM}$
would have a window of that length occurring in $t$, a contradiction.
Therefore the infinite component contributes no positive-period fixed point.

\section{All-period counts}
The component $C_\ell$ contributes all $\ell$ of its points to
$\operatorname{Fix}(T^n)$ exactly when $\ell\mid n$.  Therefore
\begin{equation}
\#\operatorname{Fix}(T^n)=\sum_{\substack{\ell\in E\\\ell\mid n}}\ell .\tag{1}
\end{equation}
Each component is a single cycle of its displayed least period.  Equivalently,
if $P(n)$ counts exact-period points, then
\[
P(n)=\sum_{d\mid n}\mu(n/d)\#\operatorname{Fix}(T^d),
\qquad C(n)=P(n)/n,
\]
so $C(1)=C(2)=C(3)=C(5)=1$ and all other primitive counts vanish.  With
\[
\zeta_Y(z)=\exp\left(\sum_{n\geq1}\#\operatorname{Fix}(T^n)\frac{z^n}{n}\right),
\]
we obtain
\begin{align*}
\log\zeta_Y(z)
 &=\sum_{\ell\in E}\sum_{q\geq1}\frac{z^{\ell q}}q
 =-\sum_{\ell\in E}\log(1-z^\ell),\\
\zeta_Y(z)&=\frac1{(1-z)(1-z^2)(1-z^3)(1-z^5)}.\tag{2}
\end{align*}

\section{Minimality obstruction and validation}
With the componentwise topology, a finite disjoint union of compact spaces is
compact.  Every tagged cycle is a proper nonempty closed invariant subset of
$Y$, so $Y$ is not minimal.  The same proof applies to any nonempty finite
union of periodic cycles disjointly attached to any nonempty aperiodic
component.  In particular, no almost-minimal conclusion is asserted.

A 60-period exact ledger and degree-30 zeta series agree
with (1)--(2).  The independent checker passes 395 assertions, SymPy passes
277 checks, replay is byte-identical, and 42 hostile cases are rejected.

The strict tuple is $(\texttt{A1\_FAIL},\texttt{A2\_FAIL},
\texttt{A3\_FAIL},\texttt{A4\_FAIL})$.  The cycles are declared topological
disjoint attachments,
not points of $X_{\rm TM}$.  We claim no target divisor, arithmetic/local
factor, target functional equation or counting law, root number, automorphy,
natural operator lift, Hilbert--P\'olya operator, or Route-B authorization.
Scope: \texttt{NO\_BAD\_EULER\_OR\_ROOT\_NUMBER}.
\end{document}
